% Producer route for the trace calibration and scale gallery.
\section{The trace calibration on the selected construction}
\label{sec:ex-lattice-trace-producer}

The trace asymptotic uses a one-dimensional conserved grade-four line and a
calibrated limiting sequence of such tensors. Here we prove the finite-
dimensional rank argument conditional on the interacting-center and
scale-dependent comparison receipts stated below. The free reference rank is
proved in Section~\ref{sec:ex-lattice-trace-free-rank}. Detailed
boundary-measure and all-radii arguments remain in Sections~\ref{sec:trace-spectrum}
and \ref{sec:trace-asymptotic}.

\subsection{The conserved tensor space has dimension one}

Let $K_0$ be the real finite-dimensional space of reference grade-four
stress tensors and $V_0$ its divergence space. Let
\[
 D_0:K_0\longrightarrow V_0
\]
be the reference divergence. Let $K$ be the interacting spin-kernel space,
$D:K\to V$ its fixed divergence, and, at scale $i$,
\[
 O_i:K\longrightarrow V_0
\]
the comparison map into the reference two-component table. Only $O_i$ depends on scale. The interacting-center dimension
equivalence and the scale-dependent comparison are source inputs stated
below; the elementary rank conclusion is proved conditional on them. The
free reference statement is not assumed: it is proved in
Section~\ref{sec:ex-lattice-trace-free-rank}.
Concretely, $K_0$ is the subspace of reference symmetric grade-four
tensors with zero scalar skim, and $V_0$ is the reference grade-five
vector table. The interacting $K$ is the centered grade-four kernel space,
and $D$ is the divergence rank map restricted to it. The optical map takes
values in the corresponding reference vector table.
\begin{enumerate}
\item $D_0$ is surjective and $\dim\ker D_0=1$;
\item $\dim K=\dim K_0$;
\item eventually $\ker O_i=\ker D$;
\item eventually there is a linear map $S_i:\mathbb R^2\to K$ such
      that $O_iS_i:\mathbb R^2\to V_0$ is bijective.
\end{enumerate}
Item 2 is the selected source's grade-four center-equivalence receipt. Its
coherent coordinate maps and centering projections are recorded in the
evidence file; this chapter does not infer it from the free reference rank.

\begin{lemma}[Rank-one conservation law]
\label{lem:ex-lattice-trace-rank-one}
Under these four properties, $\dim\ker D=1$.
\end{lemma}
\begin{proof}
Choose $i$ large enough that properties 3--4 hold. Since $O_iS_i$ is
surjective, $O_i$ is surjective. Rank--nullity and the kernel identity give
\[
 \dim\ker D=\dim\ker O_i=\dim K-\dim V_0.
\]
The reference surjection $D_0$ gives
\[
 \dim K_0=\dim V_0+\dim\ker D_0=\dim V_0+1.
\]
Property 2 identifies $\dim K=\dim K_0$, so subtraction yields
$\dim\ker D=1$.
\end{proof}

The free reference computation is proved in
Section~\ref{sec:ex-lattice-trace-free-rank}. In the notation there, the
scalar skim is the real one-point average of the $00$ tensor entry, the
free symmetric grade-four table has dimension four, and the Maxwell-reduced
grade-five vector table has dimension two. The zero-skim divergence map is
onto and its kernel is the Weyl line. In particular, this is a proved
finite-dimensional input to item 1, not a free-rank hypothesis. The source
map $T$ from the four coefficient parameters is identified with that
chapter's free stress coefficient map. Its point-probe identities and the
low-grade signed-covariance classification are proved there.


To prove item 2, let $S$ be the reference table and $T$ the interacting
table. The coherent field equivalence $U:S.Field_4\to T.Field_4$ preserves
conjugation and signed action and maps the reference unit to the interacting
unit. For a stress array $B$, define entrywise
\[
 (\mathsf U B)_{ab}=U(B_{ab}),\qquad
 (\mathsf C_T B)_{ab}=B_{ab}-\mathbb E_T(B_{ab})\,1_T.
\]
Coherence makes $\mathsf U$ preserve symmetry, reality, and signed
covariance. Centering then maps these arrays into the interacting
centered kernel space. Thus
\[
 \mathsf S(B)=\mathsf C_T(\mathsf U B)
\]
is a linear map from the reference centered kernel to the interacting one.
Define the reverse map $\mathsf R(B)=\mathsf C_S(\mathsf U^{-1}B)$.
For a centered reference array $B$, write $\lambda_{ab}=\mathbb E_T(U(B_{ab}))$.
Since $U(1_S)=1_T$,
\[
 \mathsf U^{-1}(\mathsf C_T(\mathsf U B))_{ab}
   = B_{ab}-\lambda_{ab}1_S.
\]
The reference entry is centered, so $\mathbb E_S(B_{ab})=0$; also
$\mathbb E_S(1_S)=1$. Applying $\mathsf C_S$ gives
$B_{ab}-\lambda_{ab}1_S-(-\lambda_{ab})1_S=B_{ab}$. This proves
$\mathsf R\mathsf S=I$. The same calculation with $S,T$ reversed proves
$\mathsf S\mathsf R=I$. Hence the maps are inverse linear equivalences,
which establishes item 2. This proof uses the source coherence, unit, and
centering identities; the evidence crosswalk gives their exact declarations.

Property 3 follows from the actual optical map. At grade five the source
formula is
\[
 (O_iT)(a)=b_i\,U_{i,5}^{-1}\bigl((DT)(a)\bigr),
 \qquad b_i>0,
\]
where $U_{i,5}$ is an isomorphism and
$b_i=\operatorname{crown}(Z,J,\operatorname{Stern}(V),5,i)>0$ is the
positive scalar normalization. The optical map is this scalar times
$U_{i,5}^{-1}$ composed with the divergence rank map. Therefore $O_iT=0$
exactly when $DT=0$. Positivity of $b_i$ and invertibility of $U_{i,5}$
hold at every scale $i$, in particular at all sufficiently large scales. Property 4, by contrast,
also requires the scale-dependent two-direction limits and stable inversion
established below; the free reference calculation does not prove it.

In particular, for every tensor $T$ and every scale $i$,
\begin{equation}
 O_i T=0\quad\Longleftrightarrow\quad DT=0.
 \label{eq:ex-lattice-trace-kernel-comparison}
\end{equation}
Property 4 follows once the actual source maps $S_i$ have the stated
columns $u_i^0,u_i^1$ and $\mathcal B_i=O_i\circ S_i$ converges to an
isomorphism. The stable-inversion argument is written next; the scale-
dependent convergence estimates remain source inputs.

\subsection{Solving the two-direction correction}

At scale $i$, let $u_i^0,u_i^1\in K$ be the two correction tensors and
define $S_i:\mathbb R^2\to K$ by $S_i(e_0)=u_i^0$ and $S_i(e_1)=u_i^1$.
Set $\mathcal B_i=O_i\circ S_i:\mathbb R^2\to V_0$. Let $T_i^0\in K$
be the uncorrected scaled tensor and $X_i=O_i(T_i^0)\in V_0$. Thus exact
cancellation $\mathcal B_i c_i+X_i=0$ is precisely the optical conservation
condition $O_i(T_i^0+S_i c_i)=0$. The source's limiting comparison matrix
$\mathcal B_\infty$ is an isomorphism. The following elementary lemma
isolates the only finite-dimensional inverse argument.

\begin{lemma}[Stable inversion in two directions]
\label{lem:ex-lattice-trace-stable-inversion}
Let $E,F$ be real normed spaces of dimension two. Suppose linear maps
$\mathcal B_i:E\to F$ converge in operator norm to a bijection
$\mathcal B_\infty$. Then there are $i_0$ and $C_0<\infty$ such that for
$i\geq i_0$, $\mathcal B_i$ is bijective and
\[
 \|\mathcal B_i^{-1}y\|\leq C_0\|y\|\qquad(y\in F).
\]
If $\|X_i\|\leq C_1 A_i$ with $A_i\geq0$, the unique solution
$c_i=-\mathcal B_i^{-1}X_i$ satisfies
\[
 \mathcal B_i c_i+X_i=0,
 \qquad \|c_i\|\leq C_0C_1A_i.
\]
\end{lemma}
\begin{proof}
The unit sphere of $E$ is compact. Since $\mathcal B_\infty$ is injective,
the continuous function $u\mapsto\|\mathcal B_\infty u\|$ has a positive
minimum $m$ on that sphere. For large $i$,
$\|\mathcal B_i-\mathcal B_\infty\|<m/2$, and hence
$\|\mathcal B_i u\|\geq(m/2)\|u\|$ for every $u\in E$. Thus $\mathcal
B_i$ is injective, and equality of the finite dimensions makes it
surjective. The same lower bound gives
$\|\mathcal B_i^{-1}y\|\leq(2/m)\|y\|$. Taking $C_0=2/m$ proves the
claims, including the equation and bound for $c_i$.
\end{proof}

The receiver argument is the already integrated
Lemma~\ref{lem:ex-lattice-trace-weighted-price} in
Section~\ref{sec:ex-lattice-trace-weighted-estimates}; it is not duplicated here.


We recall how the receiver lemma is applied; the symbols in this paragraph
($H$, $x$, $\rho$, $w_i$, $u_i$, $U$) are those of
Lemma~\ref{lem:ex-lattice-trace-weighted-price}, whose space $E$ we write
$\mathcal E$ to avoid a clash with the two-dimensional space $E$ of the
stable-inversion lemma.  These symbols are not used elsewhere in this
section.  In the construction, $\mathcal E$ is the real vector space
of grade-five coefficient arrays satisfying reality and signed covariance,
not the full unconstrained label space. The space $H$ is the reference
vector table. Its canonical
field map is the quotient projection applied to the grade-five coefficient
array; its kernel is the line generated by the nonzero equation-of-motion
array $x$. The scalar functional $\rho$ is the free divergence price after
that same quotient projection. The exceptional vector is the medium
\[
 w_i=\beta_i^{-1}\sigma_i^X-\sigma_i^L,
\]
where $\sigma_i^X$ is the transported chosen response, $\sigma_i^L$ is the
transported reference response at the same scale, and $\beta_i$ is the
positive scale parameter. The source proves $w_i\to x$, weighted decay of
its field image, and weighted decay under $\rho$. The finite chart gives
the eventual right inverse $u_i$, while the limiting chart gives the
section $U$. The lemma shows how these receipts produce a uniformly bounded
receiver with small weighted error; it does not prove the response, price,
or field-decay estimates themselves.

In the trace construction, $E=\mathbb R^2$ labels the two correction sectors,
$F=V_0$ is the reference table space, and $A_i$ is the reciprocal scale
price. The limiting matrix is invertible because the two reference
directions are linearly independent. Its convergence is pointwise on a
finite basis, hence in operator norm. The uncorrected divergence has norm at
most $C_1A_i$. Solving the system therefore produces coefficients with the
same $A_i$ bound and exact cancellation. Applying the flux functional to this
identity gives the coefficient normalization. The precise weighted remainder
estimates and the coefficient bound below show that the false-sector term
vanishes, while the true-sector term is $-4c_i^1+o(1)$. Since the
uncorrected term tends to $-\kappa$, the true-sector coefficient tends to
$-\kappa/4$. In the selected construction $\kappa$ is the slope constant.
The source's limiting physical tensor is therefore fixed by its Ward flux,
not by an arbitrary rescaling.

The coefficient argument uses the following estimates from the actual
selected source, stated here in mathematical form. For the two correction
directions $u_i^0,u_i^1$, a linear flux functional $\ell_i$, correction size
$A_i$, and positive weight $a_i$:
\begin{align}
 &\|X_i\|\leq C_1A_i,\qquad \|c_i\|\leq C_0C_1A_i,
   \qquad A_i\leq a_i,
   \label{eq:ex-lattice-trace-correction-bound}\\
 &\ell_i(X_i)\longrightarrow-\kappa,
   \qquad a_i\ell_i(O_iu_i^0)\longrightarrow0,
   \qquad a_i\bigl(\ell_i(O_iu_i^1)+4\bigr)\longrightarrow0,
   \label{eq:ex-lattice-trace-flux-input}\\
 &\|c_i\|\leq C_3 a_i\quad\text{eventually}.
   \label{eq:ex-lattice-trace-weighted-coefficients}
\end{align}
Here $\ell_i$ is a linear functional on $V_0$, and $\ell_i(u_i^a)$ is
shorthand for $\ell_i(O_iu_i^a)$.  The exact cancellation identity
$\mathcal B_ic_i+X_i=0$ reads, in $V_0$,
\[
 X_i+c_i^0\,O_iu_i^0+c_i^1\,O_iu_i^1=0.
\]
Since $A_i\le a_i$, the second bound in
\eqref{eq:ex-lattice-trace-correction-bound} already gives
\eqref{eq:ex-lattice-trace-weighted-coefficients} with $C_3=C_0C_1$.
By \eqref{eq:ex-lattice-trace-weighted-coefficients} and the weighted limits,
$c_i^0\ell_i(u_i^0)\to0$ and
$c_i^1(\ell_i(u_i^1)+4)\to0$. Hence
$c_i^1\ell_i(u_i^1)=-4c_i^1+o(1)$. Apply $\ell_i$ to the exact cancellation
identity. Since $\ell_i(X_i)\to-\kappa$, it gives
$-\kappa-4c_i^1+o(1)=0$, so $c_i^1\to-\kappa/4$. The selected local
estimates establish the displayed bounds; the correspondence table records
the exact producer declarations. This argument is the finite-dimensional
calibration used by the source's corrected stress sequence.

For clarity, the two limits used after solving the correction system follow
from the following explicit scale-dependent receipts. Let $B_i>0$, let
$L_i,\Omega_i,\Delta_i$ be tensors, and define
$\mathcal Y_i(T)=B_i^{-1}\mathcal P_i(T)$ for a linear scaled tensor map
$\mathcal P_i$. Put $U_{L,i}=\mathcal P_i(L_i)$,
$U_{\Omega,i}=\mathcal P_i(\Omega_i)$, and
$U_{\Delta,i}=\mathcal P_i(\Delta_i)$. Assume
\begin{align}
 U_{L,i}&\longrightarrow R_0,
 &U_{\Omega,i}&\longrightarrow R_\Omega,
 &U_{\Delta,i}&\longrightarrow R_\Delta,
 &\frac{RQ_i}{B_i}&\longrightarrow0,
 \label{eq:ex-lattice-trace-year-inputs}
\end{align}
and $T_i=B_iL_i+c_i^0\Omega_i+c_i^1\Delta_i$ with
$|c_i^0|+|c_i^1|\le C RQ_i$. For the trace limit, assume in addition
\begin{align}
 B_i\operatorname{tr}U_{L,i}&\longrightarrow0,
 &h_i\operatorname{tr}U_{\Omega,i}&\longrightarrow0,
 &|c_i^0|&\le C' h_i,
 &\operatorname{tr}U_{\Delta,i}&\longrightarrow4O,
 &c_i^1&\longrightarrow-\kappa/4.
 \label{eq:ex-lattice-trace-trace-inputs}
\end{align}
Here $h_i\ge0$ is the weight in the false-sector trace estimate. These
hypotheses keep the inverse-normalization scale $B_i$ distinct from that
weight; the source uses distinct quantities for them.
Here $O$ denotes the reference trace field, so the limit $4O$ is
field-valued.

\begin{lemma}[Normalized and trace limits from the correction receipts]
\label{lem:ex-lattice-trace-correction-limits}
Under \eqref{eq:ex-lattice-trace-year-inputs} and
\eqref{eq:ex-lattice-trace-trace-inputs},
\[
 \mathcal Y_i(T_i)\longrightarrow R_0,
 \qquad
 B_i\operatorname{tr}\mathcal Y_i(T_i)\longrightarrow-\kappa O.
\]
\end{lemma}
\begin{proof}
By linearity and the definition of $\mathcal Y_i$,
\[
 \mathcal Y_i(T_i)=U_{L,i}
 +\frac{c_i^0}{B_i}U_{\Omega,i}
 +\frac{c_i^1}{B_i}U_{\Delta,i}.
\]
The coefficient bound and $RQ_i/B_i\to0$ imply
$c_i^a/B_i\to0$ for $a=0,1$. Each $U$-sequence converges and is therefore
bounded. Both correction terms vanish, proving the first limit. Taking traces
before normalization gives
\[
 B_i\operatorname{tr}\mathcal Y_i(T_i)
 =B_i\operatorname{tr}U_{L,i}
 +c_i^0\operatorname{tr}U_{\Omega,i}
 +c_i^1\operatorname{tr}U_{\Delta,i}.
\]
The first term tends to zero. The second is bounded by
$C'h_i|\operatorname{tr}U_{\Omega,i}|$, which tends to zero. The last tends
to $(-\kappa/4)(4O)=-\kappa O$.
\end{proof}

In the source application, $\mathcal P_i$ is the scaled comparison map,
$B_i$ is the trace amplitude, and $\mathcal Y_i$ is the normalized year
map. The source's base and correction beam limits, washout ratio, base trace
vanishing, false-sector weighted trace vanishing, and true-sector trace limit
instantiate the displayed receipts. The lemma proves the normalized tensor
and trace conclusions once these estimates are installed; it does not
identify the two weights or infer one from the other.

\subsection{The calibrated stress sequence}

Throughout this section the index $i$ runs along one comparison curve: the
radii are $r_i=s_{n_i}$ and the charts are the corrected maps
$U_{i,d}=C_{i,d}$ of Theorem~\ref{thm:ex-lattice-selected-trace-gallery}.
This is the curve used in Section~\ref{sec:actual-stress-normalization}.
Write $U_i=U_{i,4}$ for its grade-four chart and, for a tensor $T$, set
\[
 Y_i(T)=r_i^4U_i^{-1}T.
\]
Apply Lemma~\ref{lem:ex-lattice-trace-correction-limits} with
$\mathcal P_i=Y_i$, so that $\mathcal Y_i=B_i^{-1}Y_i$, and put
$\widetilde T_i=B_i^{-1}T_i$ for the normalized corrected tensor.  It is
conserved, since $T_i$ is, by the comparison identity
\eqref{eq:ex-lattice-trace-kernel-comparison}.  The lemma gives
\begin{equation}
 Y_i(\widetilde T_i)\longrightarrow R_0,
 \qquad B_i\operatorname{tr}Y_i(\widetilde T_i)\longrightarrow-\kappa O,
 \label{eq:ex-lattice-trace-trial-stress}
\end{equation}
where $R_0$ is the reference stress tensor and $O$ the reference trace field.
For the physical centered tensor $R$, a separated shell-flux test on
tensors, which we denote $\phi_i$ to distinguish it from the functional
$\ell_i$ on $V_0$, has limits $\phi_i(R)\to y\neq0$ and
$\phi_i(\widetilde T_i)\to y$. Since the conserved space has dimension one,
eventually $\widetilde T_i=t_iR$. Linearity gives
$t_i=\phi_i(\widetilde T_i)/\phi_i(R)\to1$, so dividing
\eqref{eq:ex-lattice-trace-trial-stress} by $t_i$ proves
\[
 Y_i(R)\longrightarrow R_0,
 \qquad B_i\operatorname{tr}Y_i(R)\longrightarrow-\kappa O.
\]
This is precisely the normalization argument used in
Lemma~\ref{lem:trace-normalized-stress}. The source supplies the two flux
limits by its physical and reference translation Ward identities and the
separated-jet comparison. The argument above makes explicit why the common
nonzero Ward value forces the multiplicative constant to be one.

\subsection{Every vanishing positive scale has a strict subsequence}

The scale part of the source proves a slightly stronger assertion than a
single convenient scaling. Fix any sequence $s_n>0$ with $s_n\to0$. At each
integer level $d$, the finite-volume growth and funding estimates give the
following selection property: for every lower index $N$ and every $\theta>0$
there are $n>N$, a level $j\ge d$, and a finite-level chart $V_{d,j}$ such
that
\begin{equation}
 s_n<\theta,\qquad
 V_{d,j}\text{ has cutoff-}d\text{ matching receipts at location }j,
 \qquad r(V_{d,j})=s_n.
 \label{eq:ex-lattice-trace-one-stage}
\end{equation}
The matching receipts cover grades through $d$, and their accuracy is the
one requested at cutoff $d$. The estimate is uniform in the requested lower index: its
only use of $N$ is to ask the scale sequence to pass below a positive
threshold after $N$.

\begin{lemma}[Recursive strict scale selection]
\label{lem:ex-lattice-trace-strict-subsequence}
Suppose the one-stage selection property
\eqref{eq:ex-lattice-trace-one-stage} holds at every level. Then there is a
strictly increasing sequence of indices $n_d$ and matching charts $V_{d,j_d}$
at levels $j_d\ge d$ such that $s_{n_d}$ is strictly decreasing to zero and
$r(V_{d,j_d})=s_{n_d}$.
\end{lemma}
\begin{proof}
Choose an admissible triple $(n_0,j_0,V_{0,j_0})$ at level zero, using the
one-stage property with lower index zero and any positive
threshold. Having chosen $(n_d,j_d,V_{d,j_d})$, apply the property at level
$d+1$ with
lower index $n_d$ and threshold $s_{n_d}>0$. It gives
$(n_{d+1},j_{d+1},V_{d+1,j_{d+1}})$
with $n_{d+1}>n_d$ and $s_{n_{d+1}}<s_{n_d}$. Induction gives strict
monotonicity of the indices and strict decrease of the selected scales.
Since $n_d\to\infty$ and $s_n\to0$, the subsequence $s_{n_d}$ tends to
zero. The selected chart receipts hold at each level by construction.
\end{proof}

The selected construction obtains the one-stage property from its funded
finite-volume selection estimate. It finds a scale below $\theta$ past any
prescribed index, then chooses a chart with scale-location index at least
$d$ and with the accuracy requested at cutoff $d$. The recursive
construction above is the complete subsequence argument;
it does not assume that a preselected scaling already follows the requested
sequence.

\subsection{Corrected lifts and compatibility across grades}
\label{subsec:ex-trace-corrected-lifts}

The gradewise maps obtained by finite-volume extraction and projection need
not preserve the filtration inclusions. In the source, extraction and
projection are performed first; the correction is then made directly in the
selected table fields. Write $E_d$ for the reference field space in grade
$d$ and $S_d$ for that selected table field space. Let
\[
 \iota^E_d:E_d\longrightarrow E_{d+1},\qquad
 \iota^S_d:S_d\longrightarrow S_{d+1}
\]
be their inclusion maps. For $d<e$, write $\iota^E_{d,e}$ and
$\iota^S_{d,e}$ for their composites, and use the identity maps when $d=e$.
They obey the transitivity identities. The reference inclusion is injective. Choose a linear left
inverse $R_d:E_{d+1}\to E_d$, so $R_d\iota^E_d=1$. When fields carry the
finite signed symmetry and real structure, this retraction can be chosen to
preserve them: begin with any left inverse and average its conjugates over
the finite signed group. Then average that equivariant map with its
conjugate-linear reflection. The signed actions commute with the real
involution, so this second averaging preserves signed equivariance. Each
conjugate remains a left inverse because the inclusion intertwines both
actions. This is the finite-dimensional splitting used below.

For each finite-volume index $k$, let $U_{k,d}:E_d\to S_d$ be the raw chart
maps. Define the corrected maps recursively by
\begin{align}
 C_{k,0}&=U_{k,0},\\
 C_{k,d+1}&=\iota^S_d C_{k,d}R_d
       +U_{k,d+1}(1-\iota^E_dR_d).
 \label{eq:ex-trace-corrected-lift}
\end{align}
The first summand corrects the part of a field already represented at the
previous grade; the second retains the new complement. Directly using
$R_d\iota^E_d=1$ gives
\begin{align*}
 C_{k,d+1}\iota^E_d
 &=\iota^S_d C_{k,d}R_d\iota^E_d
   +U_{k,d+1}(\iota^E_d-\iota^E_dR_d\iota^E_d)\\
 &=\iota^S_d C_{k,d}.
\end{align*}
Induction and transitivity therefore give, for every pair $d\le e$,
\begin{equation}
 C_{k,e}\iota^E_{d,e}=\iota^S_{d,e}C_{k,d}.
 \label{eq:ex-trace-corrected-balanced}
\end{equation}
This is an exact identity at each index $k$, before the $k\to\infty$
limit. It is the needed balanced inclusion relation; unrelated
isomorphisms at early indices are not used to obtain it.

The same correction preserves the weak symmetry conditions by induction,
since $R_d$, both inclusions, and $U_{k,d}$ intertwine the signed and real
structures. The unit in each grade is the iterated inclusion of the
constant unit in grade zero. Thus the exact grade-zero unit identity, together
with the balanced relation, gives the unit identity at every grade. It also
preserves the tested approximation. To see the exact
finite-slot identity, let $I$ be a subset of the tuple's slots already using
$C$, and let $\mathcal W_k^I(v)$ be the tested chart correlation with $C$ in
those slots and $U$ in the others. For $i\notin I$ with $v_i\in E_{d+1}$,
put $p_i=R_dv_i$ and $q_i=(1-\iota^E_dR_d)v_i$. Linearity in that slot,
formula \eqref{eq:ex-trace-corrected-lift}, and the table inclusion identity
give the exact equality
\begin{align*}
\mathcal W_k^{I\cup\{i\}}(v)-\mathcal W_0(v)
={}&\bigl(\mathcal W_k^{I\cup\{i\}}(v[i\mapsto p_i])
             -\mathcal W_0(v[i\mapsto p_i])\bigr)\\
 &+\bigl(\mathcal W_k^I(v[i\mapsto q_i])
             -\mathcal W_0(v[i\mapsto q_i])\bigr).
\end{align*}
Here $v[i\mapsto p_i]$ has grade $d$ in slot $i$, and
$v[i\mapsto q_i]$ has grade $d+1$. Induction on the number of corrected
slots, using this grade reduction for a corrected slot, expresses the fully
corrected discrepancy for any fixed tuple as a finite sum of raw tested
discrepancies at fixed labels. Every summand tends to zero by raw approach,
so corrected correlations approach the reference. The finite sum's size and
label norms may depend on the tuple; this qualitative argument does not
retain the raw numerical rate without tracking those factors.

The raw maps $U_{k,d}$ below are formed by the finite-dimensional
ultrafilter limits and grade projections. Apply the correction after this
projection, exactly as in the source. The balanced identity
\eqref{eq:ex-trace-corrected-balanced} is then exact for each fixed pair of
grades. Once tested-correlation convergence makes the corrected maps
$C_{k,d}$ isomorphisms eventually at each fixed grade, the selected
isomorphisms obey the inclusion identity for each fixed pair on a tail of
indices. The exceptional tail may depend on
the pair of grades; no single tail uniform in all grades is asserted. This
is the compatible gradewise scaling conclusion used by the source.

Here are the finite-dimensional gluing details. At cutoff $k$, let
\[
 \mathscr L_k=\prod_{d=0}^{k}
 \operatorname{Hom}_{\mathbb C}(E_d,A_d),
\]
where $E_d$ is the free reference grade-$d$ field space and $A_d$ is the
finite family coefficient space at grade $d$. Write $S_d$ for the selected
table field space. This is finite-dimensional. The finite-volume lifts into
$\mathscr L_k$ are bounded along the one fixed ultrafilter: the source has
a uniform lift bound on its late set. The closure of a bounded set in a
finite-dimensional space is compact, so every ultrafilter on that late set
has a limit $Y_k\in\mathscr L_k$. This is done for each integer $k$ using
the same ultrafilter. Let $P_d:A_d\to S_d$ be the selected table's
projection. Put
\[
 U_{k,d}=P_d\circ Y_k(d):E_d\longrightarrow S_d
 \qquad(d\le k).
\]
The exact finite-volume reality, signed-reflection, and unit identities hold
on a set belonging to the ultrafilter. Their maps are continuous, so they
pass to $Y_k$ and then to $U_{k,d}$.

For fixed $k$, the calibrated-smearing argument also passes each
finite-volume short-test bound to the limit. Let $A_{k,v}(t)$ denote the
tested correlation of a fixed tuple $v=(v_1,\ldots,v_n)$ after applying
$U_{k,d_a}$ and scaling by $s_{n_k}$; let $A_{0,v}(t)$ denote its reference
correlation. For a tuple of arity $n\le k$, labels of grades at most $k$, and any smooth compactly
supported test $t$ inside the cutoff window, with $\sup|t|\le M$, the bound
is
\begin{equation}
 |A_{k,v}(t)-A_{0,v}(t)|
 \le \frac{1}{k+1}\,M\prod_{a=1}^n\|v_a\|.
 \label{eq:ex-lattice-trace-test-error}
\end{equation}
The change-of-variables identity for the dilated test converts the
finite-scale integral to the displayed tested correlation. For any
$\eta>0$, the upper-bound receipt makes the finite-scale discrepancy at most
the right side plus $\eta$ on a late set. Since the fixed ultrafilter lies
below that late filter, the limit obeys the same bound after letting
$\eta\downarrow0$. The cutoff window
radius tends to infinity with $k$, so it eventually contains the support of
each fixed test. Since $(k+1)^{-1}\to0$, the tested correlations converge to
their reference values. This is the approach property needed for the
scaling chart.

Fix a grade $d$. The reference smoothness and test-function uniqueness imply
that the linear functionals obtained by inserting a grade-$d$ field into
smooth tested correlations separate points of $E_d$. Since $E_d$ is
finite-dimensional, choose finitely many such tests
$\lambda_1,\ldots,\lambda_m$ spanning its dual. The raw approach
estimate gives convergence of their values for $U_{k,d}$, and the finite-sum
correction argument gives the same convergence for $C_{k,d}$. On a finite
basis the combined maps $\Lambda_k=(\lambda_{j,k})$ therefore converge in
operator norm to the injective reference map $\Lambda=(\lambda_j)$. The
minimum of $\|\Lambda v\|$ on the unit sphere is positive. Hence
$\Lambda_k$ is injective for all sufficiently large $k$, and since each
functional factors through $C_{k,d}$, so is $C_{k,d}$. The dimension receipt
gives $\dim S_d\le\dim E_d$; injectivity gives the reverse inequality.
Thus $C_{k,d}$ is an isomorphism eventually for each fixed $d$. The
exceptional tail may depend on $d$.

\begin{theorem}[Conditional gradewise ultralimit charts]
\label{thm:ex-lattice-selected-trace-gallery}
Assume the scale comparison matrices $\mathcal B_i$ converge in operator
norm to an isomorphism and $\|X_i\|\le C A_i$, so stable inversion supplies
property 4 in the rank argument; property 3 follows from
\eqref{eq:ex-lattice-trace-kernel-comparison}. Assume also the eventual
correction cancellation and estimates in
\eqref{eq:ex-lattice-trace-correction-bound}--
\eqref{eq:ex-lattice-trace-trace-inputs}, together with the one-stage
selection property \eqref{eq:ex-lattice-trace-one-stage}. For physical
stress normalization, assume separated shell-flux limits
$\phi_i(R)\to y\ne0$ and $\phi_i(\widetilde T_i)\to y$. Finally, for each cutoff
$k$, assume the finite-volume lift bound, exact unit and covariance
identities, the dimension receipt, and the short-test bound for every tuple
of arity at most $k$ described above. Then for every positive sequence
$s_n\to0$, there is a strictly increasing sequence $n_k$ such that
$s_{n_k}$ strictly decreases to zero, and there are corrected maps
$C_{k,d}:E_d\to S_d$ for $d\le k$ with the following properties:
for each fixed $d$, $C_{k,d}$ is an isomorphism for all sufficiently large
$k$; the exact unit, reality, and signed-covariance identities hold at those
indices; and for each fixed tuple and compactly supported test, the tested
correlation converges to its reference value as $k\to\infty$.  More
precisely, the uniform test estimate holds: for every fixed tuple $v$ of
reference labels, every compact set $Q$ in configuration space and every
$M>0$,
\begin{equation}\label{eq:ex-lattice-trace-uniform-tests}
 \sup_{\substack{t\in C_c^\infty,\ \operatorname{supp}t\subset Q\\
                 \|t\|_\infty\le M}}
 \bigl|A^C_{k,v}(t)-A_{0,v}(t)\bigr|\longrightarrow0,
\end{equation}
where $A^C_{k,v}$ is the tested correlation formed with the corrected maps
$C_{k,d}$. Moreover,
for each fixed pair $d\le e$, the corrected maps obey
$C_{k,e}\iota^E_{d,e}=\iota^S_{d,e}C_{k,d}$ for all sufficiently large
$k$. The tail may depend on $(d,e)$.
\end{theorem}
\begin{proof}
Apply the recursive scale-selection lemma to choose $n_k$ with strictly
decreasing $s_{n_k}$. At each cutoff $k$, the common-ultrafilter finite
dimensional limit gives the raw projected maps $U_{k,d}$ and their tested
correlation estimates. The corrected-lift calculation above produces
$C_{k,d}$; the finite slot argument transfers tested convergence, reality,
and signed covariance to these maps. The finite separating-test argument
then proves eventual injectivity grade by grade, and the dimension receipt
makes each fixed-grade map eventually bijective. The exact balanced identity
holds for every corrected map by \eqref{eq:ex-trace-corrected-balanced}; for
each fixed pair of grades, intersecting the two eventual bijectivity tails
shows the selected maps satisfy the claimed inclusion relation.

For \eqref{eq:ex-lattice-trace-uniform-tests}, the finite-slot expansion
writes $A^C_{k,v}(t)-A_{0,v}(t)$ as a finite sum of raw discrepancies
$A_{k,w}(t)-A_{0,w}(t)$ at tuples $w$ whose labels are obtained from those
of $v$ by the fixed maps $R_d$ and $1-\iota^E_dR_d$; these tuples and their
number do not depend on $k$ or on $t$.  Once the cutoff window contains
$Q$ and $k$ exceeds the arity and grades involved,
\eqref{eq:ex-lattice-trace-test-error} bounds each summand by
$(k+1)^{-1}M\prod_a\|w_a\|$.  Summing, the supremum in
\eqref{eq:ex-lattice-trace-uniform-tests} is at most $(k+1)^{-1}MK_v$ with
a constant $K_v$ depending only on $v$, and it tends to zero.\end{proof}

The corrected-lift construction proves the pairwise eventual inclusion
identities for the corrected maps; it does not claim one tail uniform over
all grades. The actual source theorem producing the full gallery also requires several
substantive analytic estimates: the one-stage
finite-volume chart selection; the scale-dependent two-direction limits;
the base and correction beam limits; the weighted false-sector flux
washout and true-sector flux convergence; the separated-jet comparison;
and the finite chart lift, unit, covariance, short-test, and dimension
receipts used in the gluing. The source crosswalk gives exact producers and
the selected public route, but this unit has not independently rederived
those estimates. Thus this chapter proves the finite-dimensional rank
calculation, the stable inversion and coefficient extraction, the scalar
subsequence recursion, and the finite-dimensional gluing conditional on
the listed source receipts; it does not claim a stand-alone proof of the
full gallery producer.
